\section{Method}

% Pipeline block diagram. Included from sections/02_method.tex.
% Requires: tikz with libraries arrows.meta, positioning, fit, backgrounds.
\begin{figure*}[t]
\centering
\resizebox{\textwidth}{!}{%
\begin{tikzpicture}[
  font=\small,
  node distance=4.5mm,
  box/.style={draw, rounded corners=1pt, align=center, minimum height=8mm,
              inner xsep=3pt, text width=20mm},
  pre/.style={box, fill=black!5},
  feat/.style={box, fill=blue!8},
  mod/.style={box, fill=green!10},
  xai/.style={box, fill=orange!12},
  diag/.style={box, fill=red!8, text width=24mm},
  ar/.style={-{Latex[length=2mm]}, thick},
  lbl/.style={font=\footnotesize\itshape}
]

% ---- stage 1: preprocessing -------------------------------------------------
\node[pre] (raw) {PCG\\ recording};
\node[pre, right=of raw] (rs) {resample\\ \SI{2}{\kilo\hertz}};
\node[pre, right=of rs] (bp) {band-pass\\ 25--\SI{400}{\hertz}\\ zero-phase};
\node[pre, right=of bp] (sp) {spike\\ removal};
\node[pre, right=of sp] (nz) {per-rec.\\ $z$-norm};
\node[pre, right=of nz] (seg) {\SI{3}{\second}\\ windows};

\draw[ar] (raw) -- (rs);
\draw[ar] (rs) -- (bp);
\draw[ar] (bp) -- (sp);
\draw[ar] (sp) -- (nz);
\draw[ar] (nz) -- (seg);

% ---- stage 2: three feature domains ----------------------------------------
\node[feat, below left=9mm and +6mm of seg] (mfcc) {MFCC\\ $13{+}\Delta{+}\Delta^2$\\ $\times 6$ stats\\ (234-d)};
\node[feat, below=9mm of seg]               (pwp)  {PWP db4\\ 12 mel bands\\ $\times 7$ desc.\\ (84-d)};
\node[feat, below right=9mm and +6mm of seg](lmel) {log-Mel\\ $32\times188$\\ (time kept)};

\draw[ar] (seg) -- (mfcc);
\draw[ar] (seg) -- (lmel);
\draw[ar] (seg) -- (pwp);

% ---- stage 3: unsupervised diagnostic --------------------------------------
\node[diag, left=12mm of mfcc] (km) {$k$-means $+$ PCA / t-SNE\\[1mm]
  ARI vs.\ \emph{diagnosis}\\ ARI vs.\ \emph{site}};
\draw[ar] (mfcc.west) -- (km.east);
\draw[ar] (pwp.south west) to[out=200,in=-80] (km.south);

% ---- stage 4: classifiers ---------------------------------------------------
\node[mod, below=15mm of mfcc] (svm) {SVM (RBF)};
\node[mod, below=15mm of lmel] (cnn) {2-D CNN\\ $\sim\!10^5$ par.};
\node[mod, below=15mm of pwp]  (rf)  {random\\ forest};

\draw[ar] (mfcc) -- (svm);
\draw[ar] (mfcc.south east) to[out=-60,in=120] (rf.north west);
\draw[ar] (pwp.south west)  to[out=-120,in=60] (svm.north east);
\draw[ar] (pwp) -- (rf);
\draw[ar] (lmel) -- (cnn);

% ---- stage 5: aggregation + evaluation -------------------------------------
\node[mod, below=9mm of cnn, text width=34mm] (agg)
  {window scores $\rightarrow$ mean probability $\rightarrow$ \textbf{recording decision}};
\draw[ar] (svm) -- ++(0,-4mm) |- (agg.west);
\draw[ar] (rf)  |- (agg.west);
\draw[ar] (cnn) -- (agg);

\node[diag, left=50mm of agg, text width=32mm] (ev)
  {MAcc $\pm$ bootstrap CI\\ McNemar $+$ Holm\\ \textbf{per-site breakdown}};
\draw[ar] (ev) -- (agg);

% ---- stage 6: XAI -----------------------------------------------------------
\node[diag, below=9mm of agg, text width=30mm] (align)
  {enrichment $E_s$ vs.\\ uniform \& shuffled null\\ $+$ sanity checks};
\node[xai, left=6mm of align, text width=26mm] (shap) {SHAP\\ (tree exact / kernel)\\ $\rightarrow$ frequency};
\node[xai, right=6mm of align, text width=26mm] (gc) {Grad-CAM\\ (last conv block)\\ $\rightarrow$ time};


\draw[ar] (svm.south) -- ++(0,-3mm) -| (shap.north);
\draw[ar] (rf.south) -- ++(0,-3mm) -| (shap.north);
\draw[ar] (cnn.south) -- ++(0,-4mm) -| (gc.north);
\draw[ar] (gc) -- (align.east|-gc.east);
\draw[ar] (shap) -- (align) ;

% ---- grouping labels --------------------------------------------------------
\begin{scope}[on background layer]
  \node[fit=(raw)(seg), draw=black!30, dashed, rounded corners,
        label={[lbl]above:preprocessing}] {};
  \node[fit=(mfcc)(lmel), draw=black!30, dashed, rounded corners,
        label={[lbl]above right:three feature domains}] {};
  \node[fit=(gc)(shap), draw=black!30, dashed, rounded corners,
        label={[lbl]below:explainability as a falsifiable test}] {};
\end{scope}

\end{tikzpicture}}
\caption{Pipeline. The three feature domains are extracted from identical
windows and evaluated under a single recording-level protocol, so that any
difference between them is attributable to the representation. The two shaded
blocks on the right are diagnostics rather than outputs: $k$-means asks whether
the class structure exists before supervision, and the enrichment test asks
whether the CNN's attribution is placed where cardiac physiology says the
evidence should be.}
\label{fig:pipeline}
\end{figure*}


Fig.~\ref{fig:pipeline} shows the pipeline. Its distinguishing property is not
any single block but the fact that the three feature domains are computed from
\emph{identical} windows and evaluated under a single recording-level protocol,
so that any difference between them is attributable to the representation and
not to the experimental setup. The remainder of this section describes each
stage and the reasoning behind the choices that are not self-evident.

\subsection{Preprocessing}

Recordings are resampled to \SI{2}{\kilo\hertz} and band-pass filtered to
25--\SI{400}{\hertz} with a fourth-order Butterworth design. Below
\SI{25}{\hertz} the signal is respiration, movement and baseline drift; above
\SI{400}{\hertz} there is essentially no cardiac acoustic energy. The band
covers S1 ($\approx$25--\SI{45}{\hertz}), S2 ($\approx$55--\SI{75}{\hertz}) and
murmur energy.

Filtering is applied forward and backward, so that it is zero-phase. This
matters specifically for the analysis in Section~\ref{sec:xai}: a causal filter
has frequency-dependent group delay, which would shift high-frequency murmur
energy relative to the low-frequency S1 fundamental, and the alignment test is
precisely a claim about \emph{when} evidence occurs.

Friction transients from the stethoscope diaphragm are suppressed by the spike
removal procedure of Schmidt et al.~\cite{schmidt2010segmentation}; these are
the loudest events in many recordings and would otherwise dominate every
energy-based descriptor. Finally each recording is $z$-normalised. Gain is a
site artifact --- different clinics used different preamplifiers --- and since
gain correlates with site and site correlates with prevalence, preserving
absolute level would hand the model a shortcut.

\subsection{Windowing}

Windows are \SI{3}{\second} long and fixed. At 60--100\,bpm a cardiac cycle
lasts 0.6--\SI{1.0}{\second}, so a \SI{3}{\second} window always contains at
least three complete cycles and is robust to the phase at which it happens to
start. Training uses \SI{50}{\percent} overlap as augmentation; validation and
test use non-overlapping windows, so that evaluation samples remain
statistically independent and confidence intervals are not artificially
narrowed.

Windows are deliberately \emph{not} cut at cardiac-cycle boundaries. Doing so
would build the location of systole into the input representation and would
make the alignment claim of Section~\ref{sec:xai} circular.

\subsection{Three feature representations}

The three representations were chosen to differ in kind rather than in
parameter setting.

\textbf{MFCC.} Thirteen cepstral coefficients from a mel filterbank, with first
and second time derivatives, each trajectory summarised over the window by six
statistics (mean, standard deviation, minimum, maximum, skewness, kurtosis),
giving 234 dimensions. This is the standard baseline in the CinC 2016
literature~\cite{davis1980comparison,rubin2016classifying}, so it anchors
comparability, and the cepstrum is a decorrelated description of spectral
shape, which is broadly what separates murmur from no-murmur.

\textbf{Log-Mel spectrogram.} The same filterbank output retained as a
two-dimensional time--frequency map and used as CNN input. This is the only
representation that preserves the time axis; the statistical pooling used for
the classical branch discards it. The contrast between the two is one of the
comparisons the study is designed to make.

\textbf{Perceptual wavelet packet (PWP).} A Daubechies-4 wavelet packet
decomposition~\cite{daubechies1988orthonormal} whose uniform subbands are merged
into perceptually spaced bands, each described by log energy, relative energy,
Shannon entropy, standard deviation, skewness, kurtosis and peak amplitude,
following the construction of Safara et al.~\cite{safara2013multi}. MFCC and
log-Mel are both STFT-based and assume local stationarity within each window,
which fits poorly a signal whose informative events are 10--\SI{50}{\milli\second}
transients; the packet transform tiles the time--frequency plane
adaptively~\cite{mallat2008wavelet,coifman1992entropy}. A further property
matters in Section~\ref{sec:xai}: PWP features map onto frequency bands exactly,
with no transform in between.

Two resolution constraints shaped these parameters and are worth stating because
they otherwise look arbitrary. First, the number of mel bands is limited by FFT
resolution: at \SI{2}{\kilo\hertz} with a 256-point FFT the bin width is
\SI{7.81}{\hertz}, so only about 48 usable bins fall inside 25--\SI{400}{\hertz},
and requesting more filters than that produces empty, all-zero filters that
inject structural zeros into every feature vector. Second, wavelet packet nodes
are not returned in frequency order: the packet tree reverses the frequency axis
in the high-pass branch at each level, so the natural ordering is a bit-reversed
permutation of frequency order, and indexing it as though it were
frequency-ordered silently scrambles the spectral axis without raising an error.
Nodes are reordered explicitly before grouping.

\subsection{Classifiers}

Three families with different inductive biases are compared: an RBF-kernel SVM
(global, distance-based, maximum-margin); a random forest (local, axis-aligned,
ensemble, and additionally admitting an exact SHAP
computation~\cite{lundberg2020local}); and a small two-dimensional CNN on
log-Mel input, comprising four Conv--BN--ReLU--MaxPool blocks, global average
pooling and a small fully connected head, on the order of $10^5$ parameters. The
network is kept deliberately small because the dataset is small; a large
pretrained network would memorise it and would obscure the feature-domain
comparison that is the point of the study.

Class imbalance is handled by inverse-frequency weighting in all three families.
Augmentation uses circular time shift --- the phase at which a window starts is
arbitrary, so the label must be invariant to it --- and additive Gaussian noise.
Frequency masking is deliberately avoided: masking a band could delete the
murmur itself, converting an abnormal example into a mislabelled normal one.

Window scores are aggregated to a recording decision by averaging predicted
probabilities, because a clinician diagnoses a patient rather than a
three-second excerpt. Majority voting was also computed and gave results within
\num{0.004} MAcc of probability averaging.
